\chapter{مناقشة النتائج وتفسيرها}

\section*{تمهيد}
يهدف هذا الفصل إلى مناقشة النتائج التي أسفرت عنها الدراسة الميدانية وتفسيرها في ضوء الإطار
النظري والدراسات السابقة، بغية الإجابة عن السؤال المركزي: \emph{كيف يدرك المراهقون (15--17 سنة)
أسلوب معاملة والديهم المتسلط، وكيف يفسّرون أثر هذا الإدراك في بناء تقديرهم لذواتهم؟} وتسير
المناقشة وفق منطقٍ منهجي موحّد يُطبَّق على كلّ بعد من أبعاد تقدير الذات الخمسة، يقوم على أربع
خطوات متتابعة: \emph{أولًا} تلخيص ما كشفت عنه المقابلات في البعد المعني؛ \emph{وثانيًا} تفسير
النتيجة في ضوء التصوّرات النظرية المؤطِّرة (بومريند، وماكوبي ومارتن، وباربر لأسلوب المعاملة
الوالدية المتسلط؛ وروزنبرغ، وكوبر سميث، وهارتر، وماسلو لتقدير الذات)؛ \emph{وثالثًا} مقارنة
النتيجة بالدراسات السابقة (أيت مولود وإكردوشن بعلي 2018، وشكراوي 2025، ولامبورن وزملائه 1991،
والدراسة النيبالية 2023)؛ \emph{ورابعًا} استخلاص خلاصة البعد. وتنتهي المناقشة بتركيبٍ عامّ
يُجيب صراحةً عن السؤال المركزي، ثمّ بإسهامات الدراسة وحدودها واستنتاجاتها وتوصياتها وآفاقها.

% =====================================================================
\section{مناقشة البعد الشخصي لتقدير الذات}
\noindent\textbf{عرض النتيجة.}\ بيّنت المقابلات أنّ البعد الشخصي كان من أكثر أبعاد تقدير الذات
ارتباطًا بإدراك التسلط؛ فقد سجّل تسعة من أصل أحد عشر مشاركًا مستوى تأثّرٍ يتراوح بين المتوسّط
والمرتفع. وتمحورت أبرز المؤشّرات حول \emph{الانسحاب من التعبير عن الرأي} و\emph{ضعف الثقة في
الرأي الشخصي والتردّد} (بنسبة 36\% لكلٍّ منهما)، و\emph{الاعتماد على أحكام الوالدين} (27\%). وقد
عبّر عدد من المشاركين عن أنّهم توقّفوا عن محاولة شرح وجهة نظرهم بعد أن أدركوا أنّها «لا تُؤخذ بعين
الاعتبار»، بل وصل الأمر ببعضهم إلى الشكّ في آرائه حتى فيما هو متأكّد منه.

\noindent\textbf{التفسير النظري.}\ يمكن فهم هذه النتيجة في ضوء تصوّر \textbf{بومريند}، الذي يصف
الأسلوب المتسلط بارتفاع الضبط ومحدودية الحوار، وهو ما يحدّ من استقلالية الأبناء ومن مشاركتهم في
القرار. ويتقاطع ذلك مع تصوّر \textbf{روزنبرغ}، الذي يرى أنّ تقدير الذات يتكوّن من التقييم الذي
يبنيه الفرد عن نفسه انطلاقًا من خبراته مع الأشخاص ذوي الأهمّية؛ فإدراك المراهق لغياب الاستماع
لرأيه ينعكس على تقييمه لقدراته وثقته بنفسه. كما يضيء تصوّر \textbf{كوبر سميث} هذه النتيجة، إذ
يجعل الشعور بالكفاءة والقدرة على الضبط الذاتي ركنًا في تقدير الذات، وهو ما يتزعزع حين يُدرك
المراهق أنّ قراراته «مُقرَّرة سلفًا».

\noindent\textbf{المقارنة بالدراسات السابقة.}\ تتّفق هذه النتيجة مع دراسة أيت مولود وإكردوشن بعلي
(2018) ودراسة شكراوي (2025) في ربط الأساليب الوالدية الضابطة بانخفاض تقدير الذات، ومع دراسة
لامبورن وزملائه (1991) في أنّ إدراك النمط المتسلط يقترن بمستويات أدنى من تقدير الذات. أمّا
الدراسة النيبالية (2023) فتختلف جزئيًا، وهو اختلاف يُردّ إلى خصوصية السياق الثقافي كما أشار
أصحابها.

\noindent\textbf{خلاصة البعد.}\ لا يرتبط تأثّر البعد الشخصي بوجود الضبط في ذاته، وإنّما بإدراك
المراهق له بوصفه تقليلًا من قيمة رأيه وتقييدًا لاستقلاليته؛ وكلّما رسخ هذا الإدراك انخفضت ثقته
بقدرته على التعبير واتّخاذ القرار. ويؤكّد ذلك وجود حالتين (م5 وم11) طوّرتا آليات تكيّف حافظت على
تقدير ذاتٍ شخصي أفضل، مقابل حالة (م9) بلغ فيها التأثّر ذروته بإضافة مؤشّر كبت المشاعر.

% =====================================================================
\section{مناقشة البعد الأسري لتقدير الذات}
\noindent\textbf{عرض النتيجة.}\ يُعدّ البعد الأسري \emph{الأعلى تأثّرًا} عبر العيّنة كاملة، إذ
سجّل تسع حالات بمستوى مرتفع أو مرتفع جدًّا. وبرز \emph{ارتباط الشعور بالقيمة بالتقبّل العاطفي
ورضا الوالدين} (27\%) و\emph{الحساسية للنقد ولوم الذات} (36\%) مؤشّرَين مشتركَين، فيما ضخّم
\emph{الصراع الأسري} هذا التأثّر لدى حالتين (م2 وم9).

\noindent\textbf{التفسير النظري.}\ تُفسَّر النتيجة في ضوء تصوّر \textbf{ماكوبي ومارتن}، اللذين
يجمعان في النمط المتسلط بين ارتفاع الضبط وانخفاض الاستجابة الانفعالية، فتقوم العلاقة على الامتثال
لا على الحوار والتفاعل الوجداني. ويعزّز ذلك تصوّر \textbf{كوبر سميث} الذي يجعل قبول الأشخاص
المهمّين وتقديرهم شرطًا لشعور الفرد بقيمته، وهو ما يتّصل مباشرةً بالمناخ الأسري. ويضيف
\textbf{ماسلو} بُعدًا مكمّلًا، إذ يجعل الحاجة إلى التقدير والانتماء حاجة نفسية أساسية يرتبط
إشباعها بالشعور بالقيمة داخل الأسرة.

\noindent\textbf{المقارنة بالدراسات السابقة.}\ تنسجم النتيجة مع دراسة أيت مولود وإكردوشن بعلي
(2018) في أنّ إدراك انخفاض الدفء وارتفاع الضبط والعقاب يرتبط بانخفاض تقدير الذات، ومع شكراوي
(2025) في وجود علاقة بين طبيعة الأساليب الوالدية ومستوى تقدير الذات، ومع لامبورن وزملائه (1991)
في اقتران النمط المتسلط بمستويات أدنى مقارنةً بالأنماط الأخرى.

\noindent\textbf{خلاصة البعد.}\ يتحدّد تأثّر البعد الأسري بالمعنى الذي يمنحه المراهق للممارسات
الوالدية؛ فحين يفسّرها تعبيرًا عن قلّة التقدير أو ضعف الاهتمام بمشاعره ينخفض شعوره بقيمته داخل
النسق الأسري، بينما خفّف إدراك «التقبّل رغم الصرامة» هذا الأثر لدى م5 وم11.

% =====================================================================
\section{مناقشة البعد الاجتماعي لتقدير الذات}
\noindent\textbf{عرض النتيجة.}\ اتّسم تأثّر البعد الاجتماعي بالتوسّط، وتمحورت مؤشّراته حول
\emph{الحذر الاجتماعي المرتفع} (45\%) و\emph{تأثير الرقابة الوالدية على بناء العلاقات} (36\%).
وأظهرت بعض الحالات أنّ العلاقات الاجتماعية تؤدّي \emph{وظيفة تعويضية} عن الضغط الأسري (م7 وم9)،
في حين مثّل م3 الحالة الأكثر تأثّرًا اجتماعيًّا نتيجة الرقابة المكثّفة.

\noindent\textbf{التفسير النظري.}\ يضيء تصوّر \textbf{باربر} هذه النتيجة عبر مفهوم \emph{التحكّم
النفسي}، الذي لا يقتصر على توجيه السلوك بل يمتدّ إلى مشاعر الأبناء وإدراكهم لذواتهم، فيحدّ من
ثقتهم في المواقف الاجتماعية. ويكمّله تصوّر \textbf{هارتر}، التي ترى أنّ تقدير الذات الاجتماعي
يتشكّل من إدراك الفرد لمدى قبوله داخل جماعات انتمائه؛ ومن ثمّ تباينت النتائج بحسب إدراك كلّ
مشارك.

\noindent\textbf{المقارنة بالدراسات السابقة.}\ تتّفق النتيجة مع الاتّجاه العام لدراسات أيت مولود
وإكردوشن بعلي (2018) وشكراوي (2025) ولامبورن وزملائه (1991) في وجود علاقة بين إدراك الأسلوب
الوالدي وتقدير الذات، غير أنّ الدراسة الحالية أضافت أنّ هذه العلاقة تتجلّى اجتماعيًّا بصورة
متفاوتة تتوقّف على تفسير المراهق لخبرته الأسرية أكثر من توقّفها على الممارسة في ذاتها.

\noindent\textbf{خلاصة البعد.}\ كلّما ارتبطت الخبرة الأسرية لدى المراهق بالخوف من النقد أو ضعف
الحوار انعكس ذلك على ثقته في التفاعل الاجتماعي، بينما ساعدت خبرات إيجابية أخرى بعض المشاركين على
الحفاظ على علاقاتهم رغم إدراكهم للتسلط.

% =====================================================================
\section{مناقشة البعد الأكاديمي لتقدير الذات}
\noindent\textbf{عرض النتيجة.}\ كان البعد الأكاديمي \emph{الأقلّ تأثّرًا}، إذ حافظ ثمانية من أحد
عشر مشاركًا على شعورٍ جيّد بالكفاءة الدراسية رغم الضغط، بل \emph{تحوّل الضغط إلى دافع للإنجاز}
لدى بعضهم (م5 وم11). واستثناءً من ذلك ربطت م8 قيمتها الذاتية بالنجاح الدراسي ربطًا مرتفعًا جعل
هذا البعد من أكثر أبعادها تأثّرًا.

\noindent\textbf{التفسير النظري.}\ يوضّح تصوّر \textbf{بومريند} أنّ النمط المتسلط يتّسم بارتفاع
الضبط وفرض القواعد؛ غير أنّ نتائج الدراسة تبيّن أنّ استجابة المراهق لا تتوقّف على وجود الضبط بل
على المعنى الذي يمنحه له، فحين أُدرِكت المتابعة الدراسية بوصفها اهتمامًا بالمستقبل تحوّلت إلى
دافع. ويعزّز ذلك تصوّر \textbf{كوبر سميث} القائل بتعدّد مجالات تقدير الذات وإمكان احتفاظ الفرد
بصورة إيجابية في مجالٍ رغم تأثّر مجالات أخرى؛ فقد شكّل النجاح الدراسي مصدرًا للشعور بالكفاءة لدى
عددٍ من المشاركين.

\noindent\textbf{المقارنة بالدراسات السابقة.}\ تتّفق النتيجة مع الدراسات السابقة في وجود علاقة
بين إدراك الأسلوب المتسلط وتقدير الذات، لكنّها تتجاوزها؛ إذ اعتمدت تلك الدراسات المنهج الكمّي
وتعاملت مع تقدير الذات متغيّرًا عامًّا، بينما كشفت الدراسة الحالية -- عبر المنهج الكيفي -- أنّ
العلاقة ليست متماثلة في جميع الأبعاد، وأنّ البعد الأكاديمي قد يبقى محصّنًا حين يفسّر المراهق
الصرامة الدراسية بوصفها حرصًا على نجاحه لا انتقاصًا من قيمته. وهذا من أبرز ما تنفرد به الدراسة.

\noindent\textbf{خلاصة البعد.}\ لا يتحدّد أثر الأسلوب المتسلط في المجال الأكاديمي بالممارسة ذاتها،
وإنّما بتأويلها؛ فحين تُؤوَّل دعمًا للنجاح تصبح مصدر دافعية لا مصدر انتقاص لتقدير الذات.

% =====================================================================
\section{مناقشة البعد الجسمي لتقدير الذات}
\noindent\textbf{عرض النتيجة.}\ احتلّ البعد الجسمي المرتبة الثانية في ضعف التأثّر بعد البعد
الأكاديمي، مع \emph{تقبّل جسدي متوسّط} لدى أغلب المشاركين. وبرزت \emph{فجوة بين الجنسين}، إذ
أظهرت الإناث حساسية أعلى تجاه التعليقات المرتبطة بالمظهر (م6، م8، م9، م10)، فيما كانت الحالات
الذكورية الأدنى تأثّرًا.

\noindent\textbf{التفسير النظري.}\ يُفسَّر ذلك في ضوء تصوّر \textbf{هارتر}، التي تجعل تقدير الذات
الجسمي من المكوّنات الأساسية في المراهقة لما يصاحبها من تغيّرات جسمية وازدياد الاهتمام بصورة
الجسد. ويكمّله تصوّر \textbf{ماسلو} في مركزية الحاجة إلى التقدير والقبول؛ فكلّما أدرك المراهق في
تعليقات والديه معنى الرفض أو المقارنة انخفض رضاه عن مظهره، بينما حافظ الشعور بالتقبّل على صورةٍ
أكثر إيجابية.

\noindent\textbf{المقارنة بالدراسات السابقة.}\ تنسجم النتيجة مع الاتّجاه العام لدراسات أيت مولود
وإكردوشن بعلي (2018) وشكراوي (2025) ولامبورن وزملائه (1991)، غير أنّ الدراسة الحالية أبرزت -- عبر
تتبّع الخبرات الذاتية -- أنّ العلاقة لا تتجلّى بالطريقة نفسها لدى الجميع، بل تختلف تبعًا لتأويل
كلّ مشارك للتعليقات الوالدية.

\noindent\textbf{خلاصة البعد.}\ لم يرتبط تأثّر البعد الجسمي بوجود التعليقات في ذاتها، وإنّما
بالمعنى الذي منحه لها المراهق؛ فحين أُدرِكت نقدًا أو مقارنةً انخفض الرضا الجسدي، وحين لم تُؤوَّل
كذلك بقيت الصورة الجسدية إيجابية.

% =====================================================================
\section{التركيب العام والإجابة عن السؤال المركزي}
تكشف القراءة التركيبية أنّ العلاقة بين إدراك أسلوب المعاملة الوالدية المتسلط وتقدير الذات لم
تتّخذ الصورة نفسها في جميع الأبعاد: فقد كان البعدان الشخصي والأسري الأكثر ارتباطًا بهذا الإدراك،
يليهما البعد الاجتماعي بتأثّر متوسّط، ثمّ البعد الجسمي بتأثّرٍ ضعيف إلى متوسّط تتخلّله فجوة بين
الجنسين، في حين ظلّ البعد الأكاديمي الأقلّ تأثّرًا، بل تحوّل الضغط فيه أحيانًا إلى دافع للإنجاز.

وبناءً عليه، يمكن الإجابة عن السؤال المركزي بالقول: إنّ المراهقين يدركون الأسلوب المتسلط من خلال
منظومةٍ من الخبرات اليومية داخل الأسرة (أسلوب الحوار، وطريقة اتّخاذ القرار، والرقابة، والعقاب،
والمقارنة، والتقدير، والقدوة، والمناخ الأسري العام)، ولا ترتبط علاقتهم بتقدير ذواتهم بهذه
الممارسات في ذاتها، وإنّما بـ\emph{المعنى} الذي يمنحونه لها. ولذلك اختلفت العلاقة من بعد إلى آخر،
ومن مراهق إلى آخر، بل بين الإخوة داخل الأسرة الواحدة، تبعًا لاختلاف الخبرات والطريقة التي فُسِّرت
بها. وتنسجم هذه الخلاصة مع تصوّرات بومريند وماكوبي ومارتن (ارتفاع الضبط وضعف الاستجابة
الانفعالية) ومع روزنبرغ وكوبر سميث (تشكّل تقدير الذات عبر الخبرات مع الأشخاص المهمّين)، لكنّها
تتجاوزها بتأكيد أنّ \emph{الإدراك} هو المتغيّر الوسيط الحاسم.

% =====================================================================
\section{إسهامات الدراسة}
لم تقتصر الدراسة على تأكيد وجود علاقة بين المتغيّرين، بل استكشفت الكيفية التي يبني بها المراهق
هذا الإدراك وينعكس على نظرته إلى ذاته. وأتاح المنهج الكيفي الكشف عن معانٍ يصعب بلوغها بالمقاييس
الكمّية، من أبرزها أنّ إدراك التسلط يتشكّل أيضًا من غياب الحوار والشعور بعدم التقدير والمقارنة بين
الإخوة والتناقض بين أقوال الوالدين وأفعالهم. كما بيّنت أنّ العلاقة ليست متجانسة عبر الأبعاد، وأنّ
البعد الأكاديمي قد يبقى محصّنًا، وهو ما يمثّل إضافةً للبحث الكيفي في البيئة المغربية.

\section{حدود الدراسة}
تُقرأ النتائج في ضوء حدودٍ منهجية: اعتماد عيّنة قصدية من أحد عشر مراهقًا لا تسمح بالتعميم
الإحصائي (وهو ما يتّسق مع طبيعة المنهج الكيفي)، وانتماء المشاركين إلى مؤسّسة تعليمية واحدة بمدينة
تطوان يجعل النتائج مرتبطة بخصوصية هذا السياق، واعتماد المقابلة شبه الموجّهة يعني أنّ النتائج تعكس
الكيفية التي أدرك بها المشاركون خبراتهم وعبّروا عنها لا وصفًا موضوعيًّا لكلّ الممارسات الوالدية.
كما برزت معطياتٌ خارج متغيّرات الدراسة (القدوة الوالدية، والصراعات الأسرية، والمناخ الأسري،
واختلاف إدراك الإخوة) تستحقّ دراسات مستقلّة.

\section{استنتاجات الدراسة}
\begin{enumerate}[leftmargin=2em,itemsep=1pt]
\item لا يتشكّل إدراك التسلط من الضبط والعقاب وحدهما، بل من مجمل الخبرات اليومية داخل الأسرة.
\item لا تقوم العلاقة بتقدير الذات على الممارسة الوالدية في ذاتها، وإنّما على المعنى الذي يمنحه
لها المراهق، ولذلك تباينت مستويات التقدير حتى داخل الأسرة الواحدة.
\item البعدان الشخصي والأسري أكثر الأبعاد ارتباطًا بإدراك التسلط.
\item العلاقة في البعد الاجتماعي غير متماثلة: تحفّظ لدى بعضهم، ووظيفة تعويضية لدى آخرين.
\item البعد الأكاديمي الأقلّ تأثّرًا، إذ فُسِّرت الصرامة الدراسية حرصًا على النجاح.
\item يتداخل إدراك التسلط مع معطياتٍ أسرية أخرى (الصراع، والمناخ الأسري، والقدوة).
\item يؤكّد ذلك أهمّية المنهج الكيفي في الكشف عن الفروق الفردية في الإدراك.
\end{enumerate}

\section{التوصيات وآفاق البحث}
\noindent\textbf{توصيات للممارسة التربوية والأسرية:} تعزيز ثقافة الحوار داخل الأسرة وإتاحة
الفرصة للمراهق للتعبير؛ وتشجيع الوالدين على الاتّساق بين توجيهاتهم وسلوكهم لما للقدوة من أثر في
الإدراك؛ وتوعيتهم بأنّ المراهق يستجيب للمعنى الذي يمنحه للممارسة لا للممارسة وحدها؛ ودعم البرامج
الإرشادية الأسرية القائمة على الحوار والتقدير مع ضبطٍ متوازن؛ وحثّ الأخصائيين على الإنصات إلى
خبرة المراهق الذاتية.\par
\noindent\textbf{آفاق البحث المستقبلية:} دراسة القدوة الوالدية وعلاقتها بإدراك الأساليب الوالدية؛
وأثر الصراعات والمناخ الأسري في هذا الإدراك؛ والفروق في الإدراك بين الإخوة داخل الأسرة الواحدة؛
وإجراء دراسات مقارنة بين المقاربتين الكيفية والكمّية؛ وتوسيع نطاق الدراسة ليشمل مؤسّسات وبيئات
ومراحل عمرية مختلفة داخل المغرب.

\section*{خلاصة الفصل}
تناول هذا الفصل مناقشة نتائج الدراسة وتفسيرها بعدًا بعدًا في ضوء الإطار النظري والدراسات السابقة،
وخلص إلى أنّ إدراك المراهق للأسلوب المتسلط -- لا الممارسة في ذاتها -- هو المتغيّر الحاسم في فهم
علاقته بتقدير الذات، وأنّ هذه العلاقة تتفاوت عبر الأبعاد وبين المشاركين، بما مكّن من الإجابة عن
السؤال المركزي وإبراز إسهامات الدراسة وحدودها وآفاقها.
